\documentclass[a4paper,oneside,onecolumn,12pt]{LegrandOrangeBook}
%\documentclass[a4paper,12pt]{article}
\usepackage[utf8]{inputenc}
\usepackage[T1]{fontenc}
\def\magyarOptions{chapterhead=unchanged}
%\usepackage[magyar]{babel}
\usepackage{longtable}
%\usepackage{amsthm}
% \usepackage{amsmath}
% \usepackage{amsfonts}
% \usepackage{txfonts}
% \usepackage{pxfonts}
% %\usepackage{eufrak}
% \usepackage{mathptmx}
\usepackage{setspace}
% \usepackage{graphicx}
\usepackage{parskip} %paragraph között legyen spacing
% \usepackage[export]{adjustbox}%kép pozicionálás
\usepackage{minted}
\usepackage{svg}
% \usepackage{hyperref}
% \usepackage{multimedia}
% \usepackage{wrapfig}
% \usepackage{csquotes}
\usepackage{ragged2e}
\usepackage{float} % For anchoring table locations

% electronic ISBN
%\usepackage[SC5b,ISBN=000-00-000-0000-0]{ean13isbn}

% printed ISBN
\usepackage[SC5b,ISBN=000-00-000-0000-0]{ean13isbn}

\newcommand{\blankpage}{%
    \newpage
    \thispagestyle{empty}
    \mbox{}
    \newpage
}

\newcommand{\comment}[1]{{\textcolor{red}{#1}}}
\newcommand{\commentaron}[1]{{\textcolor{blue}{#1}}}

% Book information for PDF metadata, remove/comment this block if not required 
\hypersetup{
	pdftitle={Real-time stock market price data analysis using neural networks}, % Title field
	pdfauthor={Bc. Eugen Fekete}, % Author field
	pdfsubject={Diplomamunka}, % Subject field
	pdfkeywords={key1, key2, key3, key4}, % Keywords
	pdfcreator={LaTeX}, % Content creator field
}

\definecolor{ocre}{RGB}{243, 102, 25} % Define the color used for highlighting throughout the book

\chapterspaceabove{6.5cm} % Default whitespace from the top of the page to the chapter title on chapter pages
\chapterspacebelow{6.75cm} % Default amount of vertical whitespace from the top margin to the start of the text on chapter pages

\onehalfspacing
%\doublespacing
\frenchspacing

\hypersetup{
    colorlinks=true,
    linkcolor=blue,
    filecolor=magenta,      
    urlcolor=cyan,
    pdftitle={Real-time stock market price data analysis using neural networks} }
\urlstyle{same}

%% \usepackage[backend=biber, style=numeric-comp, sorting=none,]{biblatex}
%% \PassOptionsToPackage{sorting=none,language=english}{biblatex}
\ExecuteBibliographyOptions{sorting=nyt}
\addbibresource{manuscript.bib}

\begin{document}

%%% Borító
\thispagestyle{empty}
\begin{minipage}[c][\textheight][c]{\textwidth}
	{\centering
	\includegraphics[keepaspectratio,width=3cm]{SelyeBanner.png}\\
	\vskip0.5cm
	{\LARGE UNIVERZITA J. SELYEHO}\\
	\vskip0.5cm
	{\LARGE SELYE JÁNOS EGYETEM}\\
    \vskip0.5cm
	{\large Fakulta ekonómie a informatiky}\\
	\vskip0.5cm
	{\large Gazdaságtudományi és Informatikai Kar}\\
	\vfill
	{\Huge Real-time stock market price data analysis using neural networks}\\
	%\vskip2cm
	%{\Huge A ZÁRÓDOLGOZAT CÍME}\\
	%\vfill
	Diplomamunka\\
	Bc. Eugen Fekete \\
    \ISBN\\
	\hfill\the\year{}, Komárno\hfill
	}
\end{minipage}

\cleardoublepage
\begingroup
\makeatletter
\let\ps@plain\ps@empty
\begin{minipage}[c][\textheight][c]{\textwidth}
	{\centering
	{\large UNIVERZITA J. SELYEHO\\SELYE JÁNOS EGYETEM}\\
	\vskip0.5cm
	{\ NÁZOV FAKULTY\\Fakulta ekonómie a informatiky\\Gazdaságtudományi és Informatikai Kar}\\
	\vfill
	{\Large NÁZOV PRÁCE\\Analýza údajov o cenách na burze v reálnom čase pomocou neurónových sietí }\\
	\vfill
	\thispagestyle{empty}
	\begin{tabular}{ll}
		Študijný program:    & Aplikovaná informatika \\
		Tanulmányi program:  & Alkalmazott Informatika\\
		Študijný odbor:      & Informatika\\
		Tanulmányi szak:     & Informatika\\
		Školiteľ:            & László Marák, PhD.\\
		Témavezető:          & László Marák, PhD.\\
		Školiace pracovisko: & Katedra informatiky\\
		Tanszék megnevezése: & Informatikai Tanszék\\
	\end{tabular}
	\vfill
	Označenie typu práce - Diplomamunka\\
	Bc. Eugen Fekete\\
    \ISBN \\
	\hfill\the\year{}, Komárno\hfill
	}
\end{minipage}
\endgroup
{
\hspace*{-2cm}
%\includegraphics[keepaspectratio, width=17cm]{./zadanie-zp_17422.pdf}
Ide jön az aláírt témakiírás
\pagebreak
}
\tableofcontents %% opcionális
\pagebreak
\listoffigures  %% opcionális
\addcontentsline{toc}{section}{Ábrák jegyzéke}

\pagebreak

\newcommand{\chpt}[1]{\chapter*{#1}\addcontentsline{toc}{section}{#1}}
% \chapterimage{kep/header.png} % Chapter heading image

\chpt{Feladatkiírás}
% \addcontentsline{toc}{chapter}{Feladatkiírás}

A szerző egy feladatot oldott meg. Még egy kicsit hosszabb, még egy kicsit hosszabb, még egy kicsit hosszabb, még egy kicsit hosszabb, még egy kicsit hosszabb, még egy kicsit hosszabb, még egy kicsit hosszabb, még egy kicsit hosszabb, még egy kicsit hosszabb, még egy kicsit hosszabb, még egy kicsit hosszabb, még egy kicsit hosszabb, még egy kicsit hosszabb, még egy kicsit hosszabb, még egy kicsit hosszabb, még egy kicsit hosszabb, még egy kicsit hosszabb, még egy kicsit hosszabb, még egy kicsit hosszabb, még egy kicsit hosszabb, még egy kicsit hosszabb, még egy kicsit hosszabb, még egy kicsit hosszabb, még egy kicsit hosszabb, még egy kicsit hosszabb, még egy kicsit hosszabb, még egy kicsit hosszabb, még egy kicsit hosszabb, még egy kicsit hosszabb, még egy kicsit hosszabb, még egy kicsit hosszabb, még egy kicsit hosszabb, még egy kicsit hosszabb, még egy kicsit hosszabb, még egy kicsit hosszabb, még egy kicsit hosszabb, még egy kicsit hosszabb, még egy kicsit hosszabb, még egy kicsit hosszabb, még egy kicsit hosszabb, még egy kicsit hosszabb, még egy kicsit hosszabb, még egy kicsit hosszabb, még egy kicsit hosszabb, még egy kicsit hosszabb, még egy kicsit hosszabb, még egy kicsit hosszabb,  

% \chapterimage{kep/header2.png} % Chapter heading image
\chpt{Opis práce}
% \addcontentsline{toc}{chapter}{Opis práce}

Autor vyrešil úlohu, ešte trošku dlhšie, ešte trošku dlhšie, ešte trošku dlhšie, ešte trošku dlhšie, ešte trošku dlhšie, ešte trošku dlhšie, ešte trošku dlhšie, ešte trošku dlhšie, ešte trošku dlhšie, ešte trošku dlhšie, ešte trošku dlhšie, ešte trošku dlhšie, ešte trošku dlhšie, ešte trošku dlhšie, ešte trošku dlhšie, ešte trošku dlhšie, ešte trošku dlhšie, ešte trošku dlhšie, ešte trošku dlhšie, ešte trošku dlhšie, ešte trošku dlhšie, ešte trošku dlhšie, ešte trošku dlhšie, ešte trošku dlhšie, ešte trošku dlhšie, ešte trošku dlhšie, ešte trošku dlhšie, ešte trošku dlhšie, ešte trošku dlhšie, ešte trošku dlhšie, ešte trošku dlhšie, ešte trošku dlhšie, ešte trošku dlhšie, ešte trošku dlhšie, ešte trošku dlhšie, ešte trošku dlhšie, ešte trošku dlhšie, ešte trošku dlhšie, ešte trošku dlhšie, ešte trošku dlhšie, ešte trošku dlhšie, ešte trošku dlhšie, ešte trošku dlhšie, ešte trošku dlhšie
\pagebreak

%\chapterimage{kep/header3.png} % Chapter heading image
\chpt{Abstrakt}\label{sec:abstrakt}
% \addcontentsline{toc}{chapter}{Abstrakt}

ešte trošku dlhšie, ešte trošku dlhšie, ešte trošku dlhšie, ešte trošku dlhšie, ešte trošku dlhšie, ešte trošku dlhšie, ešte trošku dlhšie, ešte trošku dlhšie, ešte trošku dlhšie, ešte trošku dlhšie, ešte trošku dlhšie, ešte trošku dlhšie, ešte trošku dlhšie, ešte trošku dlhšie, ešte trošku dlhšie, ešte trošku dlhšie, ešte trošku dlhšie, ešte trošku dlhšie, ešte trošku dlhšie, ešte trošku dlhšie, ešte trošku dlhšie, ešte trošku dlhšie, ešte trošku dlhšie, ešte trošku dlhšie, ešte trošku dlhšie, ešte trošku dlhšie, ešte trošku dlhšie, ešte trošku dlhšie, ešte trošku dlhšie, ešte trošku dlhšie, ešte trošku dlhšie, ešte trošku dlhšie, ešte trošku dlhšie, ešte trošku dlhšie, ešte trošku dlhšie, ešte trošku dlhšie, ešte trošku dlhšie, ešte trošku dlhšie, ešte trošku dlhšie, ešte trošku dlhšie, ešte trošku dlhšie, ešte trošku dlhšie, ešte trošku dlhšie, ešte trošku dlhšie

\textbf{Kľúčové slová: klúč1, klúč2, klúč3, }
\pagebreak

%\chapterimage{kep/header4.png} % Chapter heading image
\chpt{Absztrakt}\label{sec:absztrakt}
% \addcontentsline{toc}{chapter}{Absztrakt}

A szerző egy feladatot oldott meg. Még egy kicsit hosszabb, még egy kicsit hosszabb, még egy kicsit hosszabb, még egy kicsit hosszabb, még egy kicsit hosszabb, még egy kicsit hosszabb, még egy kicsit hosszabb, még egy kicsit hosszabb, még egy kicsit hosszabb, még egy kicsit hosszabb, még egy kicsit hosszabb, még egy kicsit hosszabb, még egy kicsit hosszabb, még egy kicsit hosszabb, még egy kicsit hosszabb, még egy kicsit hosszabb, még egy kicsit hosszabb, még egy kicsit hosszabb, még egy kicsit hosszabb, még egy kicsit hosszabb, még egy kicsit hosszabb, még egy kicsit hosszabb, még egy kicsit hosszabb, még egy kicsit hosszabb, még egy kicsit hosszabb, még egy kicsit hosszabb, még egy kicsit hosszabb, még egy kicsit hosszabb, még egy kicsit hosszabb, még egy kicsit hosszabb, még egy kicsit hosszabb, még egy kicsit hosszabb, még egy kicsit hosszabb, még egy kicsit hosszabb, még egy kicsit hosszabb, még egy kicsit hosszabb, még egy kicsit hosszabb, még egy kicsit hosszabb, még egy kicsit hosszabb, még egy kicsit hosszabb, még egy kicsit hosszabb, még egy kicsit hosszabb, még egy kicsit hosszabb, még egy kicsit hosszabb, még egy kicsit hosszabb, még egy kicsit hosszabb, még egy kicsit hosszabb, 

\textbf{Kulcsszavak: kulcs1, kulcs2, kulcs3}

\pagebreak

%\chapterimage{kep/header.png} % Chapter heading image
\chpt{Abstract}
Little longer, Little longer, Little longer, Little longer, Little longer, Little longer, Little longer, Little longer, Little longer, Little longer, Little longer, Little longer, Little longer, Little longer, Little longer, Little longer, Little longer, Little longer, Little longer, Little longer, Little longer, Little longer, Little longer, Little longer, Little longer, Little longer, Little longer, Little longer, Little longer, Little longer, Little longer, Little longer, Little longer, Little longer, Little longer, Little longer, Little longer, Little longer, Little longer, Little longer, Little longer, Little longer, Little longer, Little longer, Little longer, Little longer, Little longer, Little longer, Little longer, Little longer, Little longer, Little longer, Little longer, Little longer, Little longer, Little longer, Little longer, 

\textbf{Keywords: key1, key2, key3}

\pagebreak

\chapter*{Introduction}
\addcontentsline{toc}{chapter}{Introduction}
\markboth{}{\sffamily\normalsize{Introduction}}
% \begin{eBox}
%     Fontos megjegyzés, vagy állítás
% \end{eBox}
% \commentaron{szürke háttérrel kicsit fura nekem}\comment{válassz tetszőleges hátteret}

% \begin{figure}
%   \begin{tBox}
%     \centering{\LARGE{\url{https://google.com/}}}
%   \end{tBox}
%   \caption{Felirat}\label{fig:link}
% \end{figure}  

Machine learning (ML) plays a pivotal role in many areas of modern sciences, whether in industry, healthcare, finance and other fields. It can be used to provide a better service for the users of a search engine, a social media site or a media service provider by learning from the behaviour of the average user, predict stock prices within a specific time interval based on company performance measures and economic data, identify the risk factors for certain health conditions derived from clinical and demographic variables, identify the characters in a handwritten address from a digitized image, and so on. \cite{TESL}\\
The main objective of ML is to find rules or patterns in data to achieve certain goals. In the financial world, for example, this might involve extracting useful information from the available data to support or automate investment activities. These activities include observing the market and placing buy or sell orders based on the conclusions drawn. \cite{MLAT}

% \chapterimage{kep/header2.png} % Chapter heading image
\chapter{Elméleti rész}
Lorem ipsum dolores\\
\url{https://shayandavoodii.github.io/OnlinePortfolioSelection.jl/stable/FL/}\\

	\section{Machine Learning}
	The more common way of making a computer do work is to execute a computer program created by a human programmer. This program contains the steps and rules that turn input data into the appropriate answers, called output data. Machine learning mixes up these steps: the machine examines the input and output data, and tries to figure out what the rules should be. A system working like this is said to be trained rather than programmed. It is during the training process that the system identifies these rules by learning the patterns and relationships in the available data. \cite{DLP}\\
	Learning can be described using the definition provided by the renowned computer scientist and machine learning researcher, Tom Michael Mitchell:
	\begin{quote}
		"A computer program is said to learn from experience E with respect to some class of tasks T and performance measure P, if its performance at tasks in T, as measured by P, improves with experience E." (Tom M. Mitchell, 1997, p. 2)\cite{ML}
	\end{quote}
	As an example, a text recognition program, a so called Optical Character Recognition (OCR) software can be presented. The main goal of such a program is to correctly recognize and convert handwritten characters into digitized text. A collection of texts written in various styles is presented to the OCR software. This collection, which the system uses to learn, is called the training set, where each instance is labeled appropriately. The actual machine learning part of the software that learns and makes predictions is called the model. In this example, the task T is to recognize handwritten characters and correctly classify them, the experience E is the training set provided for learning and the performance measure P could be the accuracy of the recognition.\\
	The example mentioned described a ML system performing supervised learning and solving a classification problem. We talk about supervised learning when a training set with appropriately labeled data is available for the learning process. Two other well known types of learning are unsupervised learning, where no training set with labeled data is available, and reinforcement learning, where a software agent learns rules by interacting with its environment. A classification problem is a problem where each input can be sorted into discrete number of classes. In the previous example each letter in a text can be classified as one of the letters of the alphabet. In contrast, when predicting land prices, we do not expect discrete labels as outputs, so we can't speak of classification problems. This is known as a regression problem and we expect continuous numerical values as outputs, for which a regression algorithm is used.

		\subsection{Types of learning}
		Machine learning systems can be grouped based on the type and amount of supervision received during the training. The most common types of learning are:
		\begin{itemize}[noitemsep]
			\item Supervised learning
			\item Unsupervised learning
			\item Reinforcement learning
		\end{itemize}
		\cite{HMLSKT}
			
			\subsubsection{Supervised learning}
			For supervised learning a training set is available. The training set consists of numerous observations with predefined inputs and a corresponding output. The inputs are called predictors (also called features) and the output is referred to as the response. The observations represents individual data instances. \cite{TESL}\\
			As an example lets say we want to predict house prices and we have the following training set:\\
			\begin{table}[H]
			\begin{center}
			\begin{tabular}{|c|c|c|c||c|}
				\hline
				\textbf{Housing type} & \textbf{Size ($m^2$)} & \textbf{Number of rooms} & \textbf{Location} & \textbf{Price (€)}\\
				\hline
				Single-family home & 250 & 4 & Bratislava & 412000\\
				\hline
				Single-family home & 390 & 6 & Komárno & 186000\\ 
				\hline
				Apartment & 55 & 2 & Trnava & 118900\\
				\hline
				Single-family home & 180 & 4 & Košice & 375000\\
				\hline
				\multicolumn{5}{|c|}{\cdots} \\
				\hline
			\end{tabular}
			\end{center}
			\caption{Example training set with arbitrary values}
			\end{table}
			In Table 1.1 each row corresponds to a single house for sale. These houses are the observations and each column (excluding the last column) of the table, which represents the attributes of the houses, is an input, or a predictor. The last column is the output or the response.\\
			The variables (be that predictor or response) can be described as either quantitative or qualitative. Quantitative variables are numerical values, like the size, number of rooms and price of a house in Table 1.1. Qualitative variables take values in one of $K$ different categories, like the type of housing in Table 1.1. Problems with a quantitative response are referred to as regression problems, while problems with qualitative a response are referred to as classification problems. \cite{AISL}\\
			For each predictor measurement vector $x_i$ (where $i = 1, ..., n$), there is a corresponding response measurement $y_i$. The goal of learning is to fit a model that links the response to the predictors, either to predict the response for future observations or to gain a deeper understanding of the relationship between the response and the predictors. \cite{AISL}\\

			\textbf{Algorithms}\\
			CONTINUE: EXPLAIN THE FITTING METHOD USED IN THE FINISHED PROGRAM OR OTHER COMMON METHODS IN DEPTH\\

			\subsubsection{Unsupervised learning}
			If we only have a set of unlabeled measurement vectors $x_i$ (where $i = 1, ..., n$) with no corresponding $y_i$ response, we talk about unsupervised learning. Since we don't have an associated response, the goal of unsupervised learning is not to make predictions, but rather to analyze and discover patterns in the available data, which makes often makes working with unsupervised learning algorithms subjective. Furthermore, the lack of a corresponding response makes evaluating the obtained results difficult, as there is no way to check our work without knowing the true answer. This method of learning is often used as part of an exploratory data analysis. For example, a researcher might analyze gene expression levels in patients with certain health conditions to identify relationships or subgroups among the genes. A webstore could perform sales analyses to identify groups of customers with similar purchase histories for targeted advertisements and service customization. Similarly, a search engine might customize search results for users based on the click histories of other users with similar search patterns. \cite{AISL} \\
			For applications like these, where groups of objects or persons are identified, clustering algorithms are used. Another frequently used method in unsupervised learning is data visualization, which transforms complex, unlabeled data into a format suitable for plotting or further analysis. A related task is dimensionality reduction, where the aim is to simplify data with minimal information loss by performing feature extraction and merging correlated features. Anomaly detection algorithms can be used to detect fraudulent actions or to remove outliers from datasets during preprocessing. \cite{HMLSKT}\\

			\textbf{Algorithms}\\
			CONTINUE: EXPLAIN COMMON ALGORITHMS MORE IN DEPTH\\

			\subsubsection{Reinforcement learning}
			In reinforcement learning, a virtual object called an agent can make observation of its environment, commit actions resulting in rewards or penalties (usually defined as negative rewards) based on the outcome. \cite{HMLSKT}\\ 
			In reinforcement learning, the goal is to figure out a strategy (called a policy) that tells the agent what actions to take in different situations, in order to get the most reward over time. The policy can either be deterministic, where the the agent always picks the same action for a given situation, or stochastic, where the action is chosen probabilistically. The agent needs to learn how actions and rewards are connected, which means it needs to explore the policy space through trial and error. Policy space refers to the set of all possible policies that an agent can adopt. \cite{RLRAS}\\
			More precisely, this policy could be a neural network, which takes observations as inputs and outputs the resulting action. Alternatively, the policy could be a genetic algorithm, where each generation consists of multiple policies and new generations of policy offsprings are created based on the genetic algorithm used. Reinforcement learning is commonly used in robotics and smart home appliances like thermostats. We can also find financial securities trading support systems using reinforcement learning. \cite{HMLSKT}\\
			\textbf{Stock Trading Bot using Deep Q-Learning} \cite{STBDQ}\\
			A good example of a common approach in algorithmic trading can be illustrated by the "Stock Trading Bot using Deep Q-Learning" project created by Prabhsimran Singh. This project demonstrates an automatic stock trading bot using Deep Q-Learning, a type of reinforcement learning algorithm utilizing a deep neural network to approximate the optimal action-value function (also known as Q-function), which determines the best action to take in a given state. \cite{HCTDRL}. At any moment, the agent observes its current state (an n-day window of stock prices, which shows the stock prices of the last n days), selects and performs a buy/hold/sell action. The next state is evaluated and the agent receives a reward accordingly, based on changes in the portfolio value. The agent then updates its parameters by minimizing the computed loss. In reinforcement learning loss refers to the differnece between the target Q-value and the predicted Q-value, with the Q-value defining the optimal state-action (quality) values \cite{HMLSKT}. \\
			TEST RESULTS FROM THE CODE COMES HERE.\\
			\textbf{A Deep Reinforcement Learning Framework for the
			Financial Portfolio Management Problem} \cite{ADRLFFPMP}\\
			Another example of the use of reinforcement learning in the financial world is presented in the paper of Zhengyao Jiang, Dixing Xu and Jinjun Liang, titled "A Deep Reinforcement Learning Framework for the
			Financial Portfolio Management Problem". The paper describes the development of a system capable of automating portfolio management. Portfolio management is defined as the action of continuous reallocation of a capital into a number of financial assets. The core of the system is the Ensemble of Identical Independent Evaluators (EIIE) topology.  In an EIIE, there are multiple evaluators (models or algorithms), following the same structure or design, but each working independently examining an asset's historical data to predict future growth. The outputs of their individual evaluations are aggregated to make a final decision. The resulting scores from this evaluation are used to determine the portfolio weights for the next trading period. These weights represent the actions the agent will take. If an asset's weight increases, it'll be bought or if it goes down, it'll be sold. The portfolio weights from the last period are fed back into the EIIE for future predictions.\\
			Three different types of evaluators were tested in this work, a Convolutional Neural Network (CNN), a Recurrent Neural Network (RNN) and a Long Short Term Memory (LSTM) architecture. The framework was evaluated on a cryptocurrency exchange market in three different time period (to ensure that the tests cover different market conditions, like fast- or slow-moving market), each lasting 30 minutes. The tests included benchmarks, such as Best Stock (selects and invests in a single asset that achieved the highest final adjusted portfolio value), Uniform Buy and Hold (divides the initial capital equally among assets and holds them without rebalancing or trading), Uniform Constant
			Rebalanced Portfolio (divides the initial capital equally among assets and constantly rebalances the portfolio). Additionally, well-known strategies, such as ”Follow-the-Winner”, ”Follow-the-Loser”, ”Pattern-Matching”, and ”Meta-Learning” were also tested.\\
			At the beginning of each period, the trading agent reallocates the fund among the assets. The trading fee for certain actions on the site is calculated as the maximum possible percentage, which is 0.25\%. theframework was tested on historical data only with ideal conditions: zero slippage (meaning trade is executed immediately when an order is placed) and zero market impact (so the value of the trade is too insignificant to influence the market). For the portfolio preallocation the 11 most-volumed assets were selected. For training, historic price data from the past two days were used. The optimal policy was obtained using gradient ascent algorithm.\\
			\textbf{Algorithms}\\
			CONTINUE: EXPLAIN COMMON ALGORITHMS MORE IN DEPTH\\


		\subsection{Neural networks}


\section{Mégegy alfejezet}

felsorolás

\begin{itemize}
    \item LLE - {\it low-level} emuláció - alacsony szintű emulátorok %- ciklus pontosság
    \item HLE - {\it high-level} emuláció - magas szintű emulátorok %- absztrakciók
\end{itemize}


\subsection{Al-alfejezet}

Még egy kicsit hosszabb, még egy kicsit hosszabb, még egy kicsit hosszabb, még egy kicsit hosszabb, még egy kicsit hosszabb, még egy kicsit hosszabb, még egy kicsit hosszabb, még egy kicsit hosszabb, még egy kicsit hosszabb, még egy kicsit hosszabb, még egy kicsit hosszabb, még egy kicsit hosszabb, még egy kicsit hosszabb, még egy kicsit hosszabb, még egy kicsit hosszabb, még egy kicsit hosszabb, még egy kicsit hosszabb, még egy kicsit hosszabb, még egy kicsit hosszabb, még egy kicsit hosszabb, még egy kicsit hosszabb, még egy kicsit hosszabb, még egy kicsit hosszabb, még egy kicsit hosszabb, még egy kicsit hosszabb, még egy kicsit hosszabb, még egy kicsit hosszabb\footnote{lábjegyzet}, még egy kicsit hosszabb, még egy kicsit hosszabb, még egy kicsit hosszabb, még egy kicsit hosszabb, még egy kicsit hosszabb, még egy kicsit hosszabb, még egy kicsit hosszabb, még egy kicsit hosszabb, még egy kicsit hosszabb, még egy kicsit hosszabb, még egy kicsit hosszabb, még egy kicsit hosszabb, még egy kicsit hosszabb, még egy kicsit hosszabb, még egy kicsit hosszabb, még egy kicsit hossz

felsorolás
\begin{itemize}
    \item minden függvény-blokk atomikus, vagyis mindig lefut az elejétől a végéig, és nem szakad meg soha
    \item a függvény-blokkban nincsenek elágazások
    \item minden függvény-blokknak van egy maximum nagysága
\end{itemize}


Monospace betűtípus \texttt{fetch}, \texttt{decode} és \texttt{execute}

\paragraph{Nevesített paragrafus.} A {\it dőlt}  idézet~\cite{Ubershaders:ARidiculous}.

Beilleszteni egy képet így kell: A képre mindig kell hivatkozni a szövegből!!!

\begin{figure}[ht]
    \centering
    \includesvg[scale=0.13,inkscapelatex=false]{./kep/utasitas-abra.svg}
	\caption{Utasítás felépítése}
	\label{fig:Utasítás felépítése}
\end{figure}

\newcommand{\CS}{C${}^{\#}$}

\subsection{Így szedjük helyesen a \CS\ nyelvet}

% \chapterimage{kep/header2.png} % Chapter heading image
\chapter{Gyakorlati rész}

Még egy kicsit hosszabb, még egy kicsit hosszabb, még egy kicsit hosszabb, még egy kicsit hosszabb, még egy kicsit hosszabb, még egy kicsit hosszabb, még egy kicsit hosszabb, még egy kicsit hosszabb, még egy kicsit hosszabb, még egy kicsit hosszabb, még egy kicsit hosszabb, még egy kicsit hosszabb, még egy kicsit hosszabb, még egy kicsit hosszabb, még egy kicsit hosszabb, még egy kicsit hosszabb, még egy kicsit hosszabb, még egy kicsit hosszabb, még egy kicsit hosszabb, még egy kicsit hosszabb, még egy kicsit hosszabb, még egy kicsit hosszabb, még egy kicsit hosszabb, még egy kicsit hosszabb, még egy kicsit hosszabb, még egy kicsit hosszabb, még egy kicsit hosszabb\footnote{lábjegyzet}, még egy kicsit hosszabb, még egy kicsit hosszabb, még egy kicsit hosszabb, még egy kicsit hosszabb, még egy kicsit hosszabb, még egy kicsit hosszabb, még egy kicsit hosszabb, még egy kicsit hosszabb, még egy kicsit hosszabb, még egy kicsit hosszabb, még egy kicsit hosszabb, még egy kicsit hosszabb, még egy kicsit hosszabb, még egy kicsit hosszabb, még egy kicsit hosszabb, még egy kicsit hossz

% \chapterimage{kep/header.png} % Chapter heading image
\chapter*{Befejezés}
\addcontentsline{toc}{chapter}{Befejezés}

Még egy kicsit hosszabb, még egy kicsit hosszabb, még egy kicsit hosszabb, még egy kicsit hosszabb, még egy kicsit hosszabb, még egy kicsit hosszabb, még egy kicsit hosszabb, még egy kicsit hosszabb, még egy kicsit hosszabb, még egy kicsit hosszabb, még egy kicsit hosszabb, még egy kicsit hosszabb, még egy kicsit hosszabb, még egy kicsit hosszabb, még egy kicsit hosszabb, még egy kicsit hosszabb, még egy kicsit hosszabb, még egy kicsit hosszabb, még egy kicsit hosszabb, még egy kicsit hosszabb, még egy kicsit hosszabb, még egy kicsit hosszabb, még egy kicsit hosszabb, még egy kicsit hosszabb, még egy kicsit hosszabb, még egy kicsit hosszabb, még egy kicsit hosszabb\footnote{lábjegyzet}, még egy kicsit hosszabb, még egy kicsit hosszabb, még egy kicsit hosszabb, még egy kicsit hosszabb, még egy kicsit hosszabb, még egy kicsit hosszabb, még egy kicsit hosszabb, még egy kicsit hosszabb, még egy kicsit hosszabb, még egy kicsit hosszabb, még egy kicsit hosszabb, még egy kicsit hosszabb, még egy kicsit hosszabb, még egy kicsit hosszabb, még egy kicsit hosszabb, még egy kicsit hossz


\pagebreak
% \chapterimage{kep/header3.png} % Chapter heading image
\chapter*{Resumé}
\addcontentsline{toc}{chapter}{Resumé}
\markboth{}{\sffamily\normalsize{Resumé}}

Autor vyrešil úlohu, ešte trošku dlhšie, ešte trošku dlhšie, ešte trošku dlhšie, ešte trošku dlhšie, ešte trošku dlhšie, ešte trošku dlhšie, ešte trošku dlhšie, ešte trošku dlhšie, ešte trošku dlhšie, ešte trošku dlhšie, ešte trošku dlhšie, ešte trošku dlhšie, ešte trošku dlhšie, ešte trošku dlhšie, ešte trošku dlhšie, ešte trošku dlhšie, ešte trošku dlhšie, ešte trošku dlhšie, ešte trošku dlhšie, ešte trošku dlhšie, ešte trošku dlhšie, ešte trošku dlhšie, ešte trošku dlhšie, ešte trošku dlhšie, ešte trošku dlhšie, ešte trošku dlhšie, ešte trošku dlhšie, ešte trošku dlhšie, ešte trošku dlhšie, ešte trošku dlhšie, ešte trošku dlhšie, ešte trošku dlhšie, ešte trošku dlhšie, ešte trošku dlhšie, ešte trošku dlhšie, ešte trošku dlhšie, ešte trošku dlhšie, ešte trošku dlhšie, ešte trošku dlhšie, ešte trošku dlhšie, ešte trošku dlhšie, ešte trošku dlhšie, ešte trošku dlhšie, ešte trošku dlhšie

% \input{gitlog.tex}

\printbibliography[title=References] % Output book bibliography entries
\addcontentsline{toc}{chapter}{References}
\pagebreak
\thispagestyle{empty}
\mbox{}
\vfill
\begin{Center}
\mbox{\vskip1cm}\EANisbn
\end{Center}

\end{document}
